% =====================================================================
\chapter{A1--A7: attacking our own design}\label{ch:hard}
% =====================================================================
A design that only its authors have looked at is not yet a design. Before finalising C1--C4 we attacked our \emph{own} first version the way a
reviewer would, assuming the strongest adversary of Chapter~\ref{ch:bg}: a captured or compromised TUAV. We found seven flaws. These are flaws of our first
draft, not of the base paper; the final design described in Chapters~\ref{ch:c1}--\ref{ch:c4} closes all of them.

\begin{center}
\small
\begin{tabularx}{\linewidth}{@{}l X X l@{}}
\toprule
& \textbf{Flaw in our first design} & \textbf{Fix in the final design} & \textbf{Found by}\\
\midrule
\textbf{A1} & one forged packet naming a UAV got that UAV evicted (lock-out) & a failed signature never evicts anybody; only valid evidence counts & review\\
\textbf{A2} & a captured TUAV could mint handover tickets for anybody & the member signs its own handover request (HO2) & review, confirmed formally\\
\textbf{A3} & a captured TUAV could frame a member at a neighbour & evictions cross cells only with checkable evidence & review\\
\textbf{A4} & a credential issued for one role could claim another (a member voting as a TUAV) & the role is bound into the certified identity through $\ep$ in $h$ & review\\
\textbf{A5} & a replayed old rekey could roll the group key back & key versions $\nu$ may only increase & review\\
\textbf{A6} & a newcomer received the key already in use (and could read recent traffic) & rekey on every arrival (RK2) & review\\
\textbf{A7} & a captured old TUAV could pose as the new TUAV in a handover & the new TUAV signs its reply (HO3) & \textbf{formal verification}\\
\bottomrule
\end{tabularx}
\end{center}

\section{Why each fix works}
\begin{description}[leftmargin=1.2em,style=nextline]
\item[A1 -- availability.] A packet with a bad signature proves nothing about the identity it names: anyone can write any pseudonym into a packet.
Penalising the \emph{named} UAV therefore lets an outsider lock out any UAV. In the final design an invalid request is simply dropped (Theorem~\ref{thm:batch}
guarantees that the valid requests in the same batch still pass).
\item[A2 -- a ticket alone is not authority.] The ticket is sealed under $K_{ab}$, which a captured old TUAV knows. Requiring the member's signature on
$(\ID_a\cat\ID_b\cat H(\Tk)\cat E)$ means a ticket admits nobody without the member's own key (Theorem~\ref{thm:ho}, property 1).
\item[A3 -- evidence, not accusation.] An eviction request from one TUAV is not believed by another cell unless it carries something anyone can verify,
such as conflicting messages signed by the accused.
\item[A4 -- role binding.] With $h=\Ho(\PID\cat\ep\cat X\cat Y)$ and $\ep=0$ reserved for TUAVs, a member's credential produces a different $h$ and
fails V4 of Definition~\ref{def:cert}.
\item[A5 -- monotone versions.] An old rekey message is a valid signature on old data; accepting it again would roll the cell back to a key the attacker
may know. Members accept a rekey only if its version is larger than the current one.
\item[A6 -- backward secrecy for newcomers.] Without a rekey on arrival a newcomer receives the key that protected the last minutes of traffic.
RK2 gives it a fresh key instead (Theorem~\ref{thm:gksecrecy}(c)).
\item[A7 -- authenticate both directions.] After A2 the member was authenticated, but the reply was not: a captured old TUAV knows $\hk$ and could answer in
$\Tb$'s name. The flaw survived our own review and was found only by the symbolic model; with $\Tb$'s signature on $(\dots\cat E\cat E_b)$ all five
properties of Theorem~\ref{thm:ho} hold.
\end{description}

\begin{plainbox}
A7 is the strongest argument for formal verification in this project: it is a real man-in-the-middle position that three careful readings missed.
\end{plainbox}
